\section{实践经验：什么有效，什么无效}

\subsection{不要过度优化}

HumanLayer 坦诚承认：完全有可能花在优化 coding agent 配置上的时间比实际用 agent 写代码的时间还多---他们经历过。他们的做法是：\textbf{偏向交付}（bias towards shipping）。只在 harness 配置确实能帮助更快交付更高质量代码的程度上投入时间。

当 agent 失败时，花时间工程化一个解决方案使其不再以同样方式失败---但不会预防性地去寻找问题来解决。

\subsection{无效的做法}

\begin{table}[H]
\centering
\caption{HumanLayer 实践中无效的做法}
\begin{tabular}{@{}p{5cm}p{8cm}@{}}
\toprule
\textbf{做法} & \textbf{问题} \\
\midrule
提前设计理想 harness & 在遇到真实失败之前做的配置往往是错误的 \\
安装大量 skills 和 MCP servers ``以防万一'' & 增加无效指令负担，挤占 context \\
每次会话后运行完整测试套件（5+ 分钟） & 输出淹没 context window \\
微观优化 sub-agents 的工具访问权限 & 导致大量工具来回切换，效果更差 \\
\bottomrule
\end{tabular}
\end{table}

\subsection{有效的做法}

\begin{importantbox}{经过验证的最佳实践}
\begin{enumerate}[leftmargin=1.5em]
    \item \textbf{从简单开始}：只在 agent 实际失败时才添加配置
    \item \textbf{设计-测试-迭代-丢弃}：HumanLayer 丢弃的 hooks 远多于实际使用的
    \item \textbf{仓库级配置分发}：将经过实战检验的配置通过仓库级别的配置文件分发给整个团队
    \item \textbf{优化迭代速度}：不要追求``第一次就对''，而是让每次迭代更快
    \item \textbf{先扩展再收缩}：先给 agent 一整套能力（如 Linear），然后在了解实际需求后精简暴露给模型的内容
\end{enumerate}
\end{importantbox}

\subsection{本章小结}

Harness engineering 的最佳策略是反应式而非预防式：从简单开始，在失败中学习，逐步添加配置。同时要有勇气丢弃不管用的东西。优化迭代速度比追求一步到位更重要。


